\chapter{Introduction}
\label{ch:introduction}

\section{Motivation}
\label{sec:motivation}

Business document processing is a fundamental task in enterprise automation, affecting industries from finance to healthcare. Organizations process millions of invoices, forms, and receipts annually, requiring extraction of structured information such as \kvps{} (e.g., ``Invoice Number: 12345'', ``Total: \$137''). Manual processing is time-consuming, error-prone, and costly, motivating the need for automated solutions.

Traditional approaches rely on template matching or rule-based systems, which lack generalization to diverse document layouts and formats. Recent advances in deep learning, particularly multimodal pre-trained models like LayoutLMv3 \cite{layoutlmv3}, offer promising solutions by jointly modeling text, layout, and visual features.

\section{Problem Statement}
\label{sec:problem}

Given a business document image, the task is to extract all \kvps{} such that three requirements are met simultaneously: the textual content of each key and value must be recognised correctly (\emph{text accuracy}); the bounding boxes around those text regions must be localised precisely (\emph{location accuracy}); and each key must be associated with its corresponding value (\emph{linking accuracy}). Several factors make this difficult. Documents vary widely in layout, ranging from sparse forms to dense invoices; keys and values may sit close together or far apart on the page; scans are frequently degraded by rotation, occlusion, or other visual noise; and the key--value structure is often implicit rather than explicitly encoded in the layout.

\section{Research Questions}
\label{sec:research_questions}

This thesis addresses the following research questions:

\begin{enumerate}
    \item \textbf{RQ1}: How do layout characteristics vary across business documents, and how can this inform model design?
    
    \item \textbf{RQ2}: Can a discriminative approach using LayoutLMv3 with a learned biaffine linker serve as a viable alternative to generative baselines for \kvp{} extraction on the KVP10k benchmark, which did not empirically evaluate such encoder-based models?
    
    \item \textbf{RQ3}: How does linking granularity (token-level vs.\ span-level) affect key-value association accuracy?
    
    \item \textbf{RQ4}: What are the key factors contributing to model success and failure cases, and how do they vary across document layout types?
\end{enumerate}

\section{Contributions}
\label{sec:contributions}

This thesis makes the following contributions:

\begin{enumerate}
    \item \textbf{Dataset Analysis}: Comprehensive characterization of KVP10k dataset including layout clustering, spatial pattern analysis, and annotation statistics (Chapter~\ref{ch:dataset})
    
    \item \textbf{Architecture}: Implementation of LayoutLMv3 + Biaffine Linker for discriminative \kvp{} extraction with explicit key-value association, in both token-level (V1) and span-level (V2/V4) variants (Chapter~\ref{ch:methodology})
    
    \item \textbf{Empirical Evaluation on KVP10k}: Systematic benchmark on the KVP10k test set using the dual-threshold evaluation protocol (NED, IoU) from Naparstek et al.\ (2024), directly testing the claim that encoder-based models are insufficient for this task (Chapter~\ref{ch:exresults})
    
    \item \textbf{Linking Granularity Analysis}: Detailed investigation of why token-level linking fails ($\sim$450:1 class imbalance) and how span-level linking addresses this (Chapter~\ref{ch:exresults})
    
    \item \textbf{Diagnostic Pipeline Analysis}: Systematic identification and correction of three compounding data-pipeline bugs, yielding regular-KVP macro F1\,=\,0.334 (text) / 0.309 (text+location) for corrected V4 (Chapter~\ref{ch:exresults})
    
    \item \textbf{Error Analysis}: Per-cluster breakdown showing layout-dependent failure modes for both generative and discriminative approaches (Chapter~\ref{ch:exresults})
\end{enumerate}

\section{Thesis Structure}
\label{sec:structure}

The remainder of this thesis is organized as follows. Chapter~\ref{ch:literature} reviews related work in document understanding, \kvp{} extraction, and multimodal pre-training. Chapter~\ref{ch:dataset} presents the KVP10k dataset together with the layout clustering analysis and spatial pattern characterization. Chapter~\ref{ch:methodology} describes the proposed approach, including the LayoutLMv3 encoder, the biaffine linker architecture, and the evaluation protocol. Chapter~\ref{ch:implementation} details the implementation, covering data preprocessing, model architecture, the training procedure, and the HPC setup. Chapter~\ref{ch:exresults} presents the experiments and results in a story-driven format, moving from the generative baseline through token-level and span-level linking and the diagnostic analysis to the corrected V4 model. Chapter~\ref{ch:discussion} discusses the findings, the generative-versus-discriminative trade-offs, limitations, and future work, and Chapter~\ref{ch:conclusion} summarizes the contributions and conclusions.
